% =====================================================================
% Appendix (separate page numbering: A, B, C)
% =====================================================================
\section*{Appendix A: Proofs and Supplementary Derivations}
\addcontentsline{toc}{section}{Appendix A: Proofs and Supplementary Derivations}

\subsection*{A.1 Problem 1: Intersection of the Optimal Arc and the Proof That C3 Adds No Occlusion}
\addcontentsline{toc}{subsection}{A.1 Problem 1: Intersection of the optimal arc and C3}

\noindent\textbf{(a) Derivation of the blocked arc.} Let $T'=(84.6410,70.0000)$ and $\rmin=7.5821$~m. The $1.5$~m expanded rectangle of B2 is $[61.5,80.5]\times[57.5,72.5]$; a point on the circle, $(84.6410+7.5821\cos\varphi,\ 70+7.5821\sin\varphi)$, falls inside this rectangle if and only if
\begin{equation}
  \cos\varphi\le\frac{80.5-84.6410}{7.5821}=-0.54616,
  \qquad
  \sin\varphi\le\frac{72.5-70}{7.5821}=0.32972.
\end{equation}
The first inequality gives $\varphi\in[123.11^\circ,236.89^\circ]$, and the second gives $\varphi\in[0^\circ,19.25^\circ]\cup[160.75^\circ,360^\circ]$ ($x\ge61.5$ and $y\ge57.5$ hold automatically on the circle because $x\ge77.06$ and $y\ge62.42$ there). Their intersection is $\varphi\in[160.75^\circ,236.89^\circ]$, i.e., a blocked arc of $76.15^\circ$. Hence C1 leaves the continuous arc $[236.89^\circ,360^\circ)\cup(0^\circ,160.75^\circ]$ on the circle, of total length $283.85^\circ$.

\noindent\textbf{(b) C3 adds no occlusion on the feasible arc.} Basis: the segment $pT'$ is contained in the disk centered at $T'$ with radius $\rmin$ (convexity of the disk); the distances from the other buildings to $T'$ are $67.09$, $60.64$, $36.00$, and $41.32$~m, all greater than $\rmin+1.5=9.0821$~m (see Section~4.2.3), so only B2 needs to be excluded. For $p$ on the feasible arc, consider two cases.

\noindent (i) $x(p)\ge79$: the $x$-coordinate of the horizontal projection of the LOS lies between $x(p)$ and $84.6410$, hence is always $\ge79$ and outside the east face of B2.

\noindent (ii) $x(p)<79$: by the C1 test in (a), $p$ must satisfy $y(p)>72.5$. Let $p=(x_0,y_0)$, $t=y_0-70\in(2.5,7.5821]$, and $x_0=84.6410-\sqrt{\rmin^2-t^2}$. The ordinate of the intersection of $pT'$ with the line $x=79$ reduces to
\begin{equation}
  y(79)=70+\frac{5.6410\,t}{\sqrt{\rmin^2-t^2}}.
\end{equation}
This expression is strictly increasing on $t\in(2.5,7.5821]$, so $y(79)\ge y(79)|_{t=2.5}=71.97>71$; that is, the LOS passes north of the north face of B2 ($y=71$) and does not enter the B2 footprint. Therefore C3 produces no additional occlusion on the feasible arc.

\subsection*{A.2 Problem 2(1): Algebraic Expansion of $\Sigma=A^{-1}DA^{-T}$}
\addcontentsline{toc}{subsection}{A.2 Problem 2(1): Algebraic expansion of the covariance}

Let the LOS azimuths of stations $p_1$ and $p_2$ be $\alpha_i$, with unit directions $u_i=(\cos\alpha_i,\sin\alpha_i)$ and unit normals $n_i=(-\sin\alpha_i,\cos\alpha_i)$ ($n_i^Tu_i=0$). A first-order azimuth error $\delta\alpha_i$ makes the signed distance from the true target to the perturbed bearing line approximately $d_i\delta\alpha_i$, so
\begin{equation}
  n_i^T\Delta p=d_i\delta\alpha_i,\qquad i=1,2,
\end{equation}
which in matrix form is $A\Delta p=\delta$ with
\begin{equation}
  A=\begin{bmatrix}n_1^T\\ n_2^T\end{bmatrix}
  =\begin{bmatrix}n_{1x}&n_{1y}\\ n_{2x}&n_{2y}\end{bmatrix},
  \qquad
  \delta=(d_1\delta\alpha_1,d_2\delta\alpha_2)^T.
\end{equation}
Since $\delta\alpha_1$ and $\delta\alpha_2$ are independent and identically distributed $\mathcal{N}(0,\sigma^2)$, $\mathbb{E}[\delta\delta^T]=D=\diag(d_1^2\sigma^2,d_2^2\sigma^2)$, and
\begin{equation}
  \Sigma=\mathbb{E}[\Delta p\Delta p^T]
  =A^{-1}\mathbb{E}[\delta\delta^T]A^{-T}=A^{-1}DA^{-T}.
\end{equation}

\noindent\textbf{(a) Inverse and columns.} $\det A=n_{1x}n_{2y}-n_{1y}n_{2x}=n_1\times n_2=\sin\gamma$ ($\gamma$ is the angle between the two LOS; the signed area is $\pm\sin\gamma$, and its absolute value enters the variance). The $2\times2$ adjugate is
\begin{equation}
  \operatorname{adj}A=\begin{bmatrix}n_{2y}&-n_{1y}\\ -n_{2x}&n_{1x}\end{bmatrix},
\end{equation}
so
\begin{equation}
  A^{-1}=\frac{1}{\det A}\begin{bmatrix}n_{2y}&-n_{1y}\\ -n_{2x}&n_{1x}\end{bmatrix}.
\end{equation}
Its two columns are
\begin{equation}
  \mathrm{col}_1(A^{-1})=\frac{1}{\det A}(n_{2y},-n_{2x})^T=\frac{1}{\det A}u_2,
\end{equation}
\begin{equation}
  \mathrm{col}_2(A^{-1})=\frac{1}{\det A}(-n_{1y},n_{1x})^T=-\frac{1}{\det A}u_1,
\end{equation}
using the $90^\circ$ rotation relation between $u_i=(\cos\alpha_i,\sin\alpha_i)$ and $n_i=(-\sin\alpha_i,\cos\alpha_i)$.

\noindent\textbf{(b) Outer-product expansion.}
\begin{equation}
  \Sigma=A^{-1}DA^{-T}
  =\sum_{i=1}^{2}D_{ii}\,\mathrm{col}_i(A^{-1})\,\mathrm{col}_i(A^{-1})^T
  =\frac{\sigma^2}{\det^2A}\left[d_1^2u_2u_2^T+d_2^2u_1u_1^T\right].
\end{equation}
No cross term appears because $D$ is diagonal; if $u_i$ were mistakenly written as the normal $n_i$, the expression would be rotated by $90^\circ$ (trace and determinant unchanged, but the matrix itself different). The main text consistently uses the LOS direction $u_i$.

\noindent\textbf{(c) Trace and determinant.}
\begin{equation}
  \trace\Sigma=\frac{\sigma^2}{\det^2A}\left[d_1^2\|u_2\|^2+d_2^2\|u_1\|^2\right]
  =\frac{\sigma^2(d_1^2+d_2^2)}{\sin^2\gamma},
\end{equation}
\begin{equation}
  \det\Sigma=\frac{\det D}{\det^2A}=\frac{\sigma^4d_1^2d_2^2}{\sin^2\gamma}.
\end{equation}
These correspond respectively to
\begin{equation}
  \mathrm{DRMS}^2=\frac{\trace\Sigma}{2}
  =\frac{\sigma^2(d_1^2+d_2^2)}{2\sin^2\gamma}
  \ \Longrightarrow\
  \mathrm{DRMS}=\frac{\sigma\sqrt{(d_1^2+d_2^2)/2}}{|\sin\gamma|}=\sigma G.
\end{equation}

\noindent\textbf{(d) $1\sigma$ ellipse and $95\%$ radius.} From $\det\Sigma$, the area of the $1\sigma$ error ellipse is
\begin{equation}
  S_{1\sigma}=\pi\sqrt{\det\Sigma}=\frac{\pi\sigma^2d_1d_2}{|\sin\gamma|}.
\end{equation}
When $d_1=d_2=d$, $S_{1\sigma}=\pi\sigma^2dG$, so the two are of the same order; in general, both $G$ and $S_{1\sigma}$ vary monotonically and inversely with $|\sin\gamma|$, so $G$ is a monotone representative of ``area-type'' accuracy indices. The $95\%$ confidence ellipse is
\begin{equation}
  (\Delta p)^T\Sigma^{-1}(\Delta p)\le\chi^2_2(0.95)=5.991,
\end{equation}
with semi-axes $\sqrt{5.991\lambda_i}$, where $\lambda_i$ are the eigenvalues of $\Sigma$. When $\Sigma$ is isotropic ($\lambda_1=\lambda_2=\mathrm{DRMS}^2$), the ellipse degenerates to a circle of radius $2.4477\,\mathrm{DRMS}$; in general, $2.4477\,\mathrm{DRMS}$ is the nominal circle radius, and exact $95\%$ coverage should use the ellipse semi-axes. The main text reports $R_{95}=2.4477\,\mathrm{DRMS}$ and notes ``exact when isotropic.''

\noindent\textbf{(e) Degeneracy and condition number.} As $\gamma\to0^\circ$ or $180^\circ$, $\det A=\pm\sin\gamma\to0$, $A$ becomes singular, and the entries of $A^{-1}$ diverge as $1/|\sin\gamma|$; by (c), $\trace\Sigma$ and $\det\Sigma$ both diverge as $1/\sin^2\gamma$. $\det A=n_1\times n_2=\pm\sin\gamma$ measures the singularity of $A$; using $1/|\det A|=1/|\sin\gamma|$ as the condition-number indicator (same order as the 2-norm condition number $\cot(\gamma/2)$, both diverging as $1/\gamma$ as $\gamma\to0$) gives the quantitative basis for the operating window $[30^\circ,150^\circ]$: at $\gamma=30^\circ$, $G=2G(90^\circ)$, i.e., the accuracy cost doubles.

\subsection*{A.3 Problem 3: Linearization Residual RMS $=\sigma(s)^2/d^2$ and Hermite Quadrature}
\addcontentsline{toc}{subsection}{A.3 Problem 3: Linearization residual and Hermite quadrature}

\noindent\textbf{(a) Exact angle and first-order approximation.} Take the true position as the origin and the true LOS as the $x'$ axis, with the target at $(d,0)$, $d=40$~m. Decompose the localization error into along-range $\delta_\parallel$ and cross-range $\delta_\perp$, isotropic and independent, $\delta_\parallel,\delta_\perp\sim\mathcal{N}(0,\sigma(s)^2)$. The exact azimuth deviation from the estimated position to the target is
\begin{equation}
  \Delta\alpha=\operatorname{atan2}(-\delta_\perp,d-\delta_\parallel).
\end{equation}
Let $u=\delta_\parallel/d$ and $v=\delta_\perp/d$; then $d-\delta_\parallel=d(1-u)$, and
\begin{equation}
  \Delta\alpha=-\arctan\frac{v}{1-u}.
\end{equation}
Expanding $v/(1-u)=v(1+u+u^2+O(u^3))$ for $|u|<1$ and using $\arctan z=z-z^3/3+O(z^5)$,
\begin{equation}
  \Delta\alpha=-\left[v+uv+u^2v-\frac{v^3}{3}+O(\|(u,v)\|^4)\right].
\end{equation}
The first-order (linearized) model takes $\Delta\alpha\approx-v=-\delta_\perp/d$.

\noindent\textbf{(b) Residual and RMS.} The residual is
\begin{equation}
  R=\Delta\alpha-(-v)=-uv-u^2v+\frac{v^3}{3}+O(\|(u,v)\|^4).
\end{equation}
The leading term is $-uv$. Since $\delta_\parallel$ and $\delta_\perp$ are independent, zero-mean, and have variance $\sigma^2$,
\begin{equation}
  \mathbb{E}[R^2]=\mathbb{E}[u^2v^2]+O(\sigma^6/d^6)
  =\frac{\mathbb{E}[\delta_\parallel^2]\,\mathbb{E}[\delta_\perp^2]}{d^4}+O(\sigma^6/d^6)
  =\frac{\sigma^4}{d^4}+O(\sigma^6/d^6).
\end{equation}
Hence the residual RMS (in radians) is
\begin{equation}
  \mathrm{RMS}_R=\sqrt{\mathbb{E}[R^2]}=\frac{\sigma(s)^2}{d^2}+O(\sigma(s)^3/d^3).
\end{equation}
Note that $\sigma^2/d^2$ is dimensionless; under the small-angle approximation its numerical value is the RMS in radians, and conversion to degrees multiplies by $180/\pi$. At $s=80$~m, $\sigma=1.9$~m, and $d=40$~m:
\begin{equation}
  \mathrm{RMS}_R=\left(\frac{1.9}{40}\right)^2=2.256\times10^{-3}~\mathrm{rad}=0.1293^\circ.
\end{equation}
Relative to $\sigma_\psi(80)=2.8995^\circ$ at the same point, the standard-deviation share is $0.1293/2.8995=4.46\%$, and the variance share is $(0.1293/2.8995)^2=0.199\%\approx0.2\%$. Higher-order terms are $O(\sigma^3/d^3)$, and at $\sigma/d\le1.9/40=0.0475$ they are at most about $5\%$ of the leading term, so the first-order residual estimate is consistent within the parameter range of this paper.

\noindent\textbf{(c) Normal composition and the $P$ criterion.} The gimbal bias $\varepsilon_0\sim\mathcal{N}(0,0.8^2)$, the control error $\varepsilon_c\sim\mathcal{N}(0,0.6^2)$, and the localization error are independent, so $S=\varepsilon_0+\varepsilon_c\sim\mathcal{N}(0,1.0^2)$ (in degrees). Given $(\delta_\parallel,\delta_\perp)$, $\Delta\alpha$ is deterministic, and the conditional probability is
\begin{equation}
  P(|\Delta\alpha+S|\le5^\circ\mid\delta_\parallel,\delta_\perp)
  =\Phi\!\left(\frac{5^\circ-\Delta\alpha}{1^\circ}\right)
  -\Phi\!\left(\frac{-5^\circ-\Delta\alpha}{1^\circ}\right),
\end{equation}
where $\Phi$ is the standard normal CDF. Therefore
\begin{equation}
  P(s)=\mathbb{E}_{\delta\parallel,\delta\perp}\!\left[
  \Phi\!\left(\frac{5^\circ-\Delta\alpha(\delta_\parallel,\delta_\perp)}{1^\circ}\right)
  -\Phi\!\left(\frac{-5^\circ-\Delta\alpha(\delta_\parallel,\delta_\perp)}{1^\circ}\right)\right].
\end{equation}
Since $\delta_\parallel$ and $\delta_\perp$ are independent and have equal variance, the two-dimensional expectation can be evaluated by tensor-product Gauss--Hermite quadrature:
\begin{equation}
  P(s)\approx\frac{1}{\pi}\sum_{i=1}^{256}\sum_{j=1}^{256}w_iw_j
  \left[\Phi\!\left(\frac{5^\circ-\Delta\alpha(\sqrt2\sigma(s)x_i,\sqrt2\sigma(s)x_j)}{1^\circ}\right)
  -\Phi\!\left(\frac{-5^\circ-\Delta\alpha(\cdot)}{1^\circ}\right)\right],
\end{equation}
where $x_i$ and $w_i$ are the 256-point probabilistic Hermite nodes and weights ($\int e^{-x^2}f(x)\,dx\approx\sum w_if(x_i)$, $\sum w_i=\sqrt\pi$). $P(s)$ decreases monotonically with $s$ ($\sigma(s)=a+bs$ is increasing, and the distribution of $\Delta\alpha$ widens with $\sigma$); bisection/root finding on $P(s)=0.95$ gives $s_{\max}=66.84$~m. The closed-form linearization gives $66.92$~m, a $0.13\%$ difference consistent with the Monte Carlo results in Table~\ref{tab:4-5}.

% ---------------------------------------------------------------------
\section*{Appendix B: Algorithm Routes and Main Program List}
\addcontentsline{toc}{section}{Appendix B: Algorithm Routes and Main Program List}

\subsection*{B.1 Algorithm routes (full table)}
\renewcommand{\thetable}{B-\arabic{table}}
\setcounter{table}{0}
\begin{table}[htbp]
\centering
\caption{Algorithm routes and selection rationale for the four problems (full)}
\label{tab:B-1}
{\footnotesize
\renewcommand{\arraystretch}{1.08}
\begin{tabularnb}{@{}p{10mm}p{43mm}p{37mm}p{26mm}p{32mm}@{}}
  \toprule
  Problem & Algorithm (key steps) & Rationale & Scale / complexity & Comparison with alternatives \\
  \midrule
  Q1 & Closed-form lower bound (C2 $\Rightarrow r\ge\rmin$); 200,000-point sampling on the $\rmin$ circle; 3-D slab intersection; visibility graph + Dijkstra; three-level lexicographic rule &
  $Q=1/r$ is strictly monotone, so the optimal value is fixed in closed form by the active constraint; the optimal set is a one-dimensional arc, and sampling only describes the set, not a 2-D search &
  Closed form $O(1)$; arc sampling $O(2\times10^5)$; one grid verification $O(10^6)$ &
  A 0.1~m grid only gives the best grid value 7.5828~m; the closed form also gives the global certificate 7.5821~m \\
  \midrule
  Q2(1) & Covariance propagation $\Sigma=A^{-1}DA^{-T}$; QM--AM lower-bound certificate; enumeration of $C(5,2)=10$ pairs; polar candidate stations + Nelder--Mead &
  The equality condition reduces optimality verification to a one-dimensional $90^\circ$ pairing test on the $\rmin$ circle; with only five vehicles, pairs can be exhaustively enumerated &
  10 pairs; 4-D NM $\le300$ iterations per pair &
  Direct 4-D search has no optimality proof and easily falls into local optima; the certificate gives $G^*=5.8188$~m \\
  \midrule
  Q2(2) & 1,000 assignment structures expanded through reuse-chain orderings into 5,880 combinations; fixed-point timeline simulation; greedy initialization $\to$ Gauss--Seidel coordinate descent $\to$ Nelder--Mead; top-12 re-refinement; deadlock detection &
  With only 5 vehicles and 6 slots, complete enumeration exhausts the structures and is deterministic and reproducible; continuous station positions are handled by local refinement &
  5,880 simulations; at most 2 tasks per vehicle; $\le160$ NM iterations per call &
  10.355~s vs.\ 20.968~s for nearest assignment ($-50.6\%$); certified LB 7.391~s (gap 28.6\%); randomized searches such as GA cannot guarantee the same solution \\
  \midrule
  Q3 & Closed-form first-order propagation; two-sided quantile criterion; bisection; $256\times256$ Gauss--Hermite quadrature; $10^6$ MC verification &
  $\sigma_\psi(s)$ is monotone, so the $P$ criterion can be inverted in closed form; Hermite quadrature has no sampling error and uses $1/15$ of the MC nodes &
  Closed form $O(1)$; 65,536 quadrature nodes; $10^6$ MC for verification only &
  Closed form 66.92~m vs.\ exact 66.84~m (0.13\%); direct MC root search requires multiple rounds of $10^6$ samples \\
  \midrule
  Q4 & Event-triggered snapshot; reuse of the Q2(2) solver; deterministic enumeration of 74 scenarios; keep/simple-repair/re-optimization comparison &
  After a disturbance, available resources and pending alerts are of the same order as Q2(2); reusing the solver keeps the conventions consistent; the scenario set is small, so no online stochastic optimization is needed; simple repair separates reaction from re-optimization &
  74 scenarios $\times$ one re-optimization each; timeline audit finds 0 negative durations/overlaps &
  Simple repair restores $\eta$ in 67/74; re-optimization lowers $t_{\mathrm{last}}$ in 61/74 (mean 4.35~s) \\
  \bottomrule
\end{tabularnb}
}
\end{table}

\subsection*{B.2 Main program list}
\begin{itemize}
\item \textbf{Problem 1} (\texttt{Q1/}): \texttt{q1\_solve.py} (domain gridding, arc sampling, tie-breaking, sensitivity); \texttt{area\_estimators.py}, \texttt{q1\_metric\_compare.py}, \texttt{export\_params.py} (three-convention area recomputation, $Q/Q_2$ comparison, export of the B1--B5/V1--V5/T1--T3 parameters).
\item \textbf{Problem 2} (\texttt{Q2/}): \texttt{q2a\_core.py}, \texttt{q2b\_core.py}, \texttt{q2b\_gamma.py}, \texttt{q2b\_bound.py} (two-station accuracy, vehicle and station selection, six-slot five-vehicle scheduling and baseline, 2(2) intersection-geometry check, reuse-aware certified lower bound).
\item \textbf{Problem 3} (\texttt{Q3/}): \texttt{q3\_solve.py} (error propagation, $s_{\max}$, the $(a,b)$ region, Monte Carlo, quadrature).
\item \textbf{Problem 4} (\texttt{Q4/}): \texttt{q4\_core.py}, \texttt{q4\_repair.py}, \texttt{q4\_robust.py}, \texttt{q4\_tlast\_diff.py}, \texttt{q4\_scenario\_list.py} (state snapshot, rescheduling, simple-repair baseline, 74 scenarios, robustness checks, timeline comparison, scenario details).
\item \textbf{Build scripts}: \texttt{build\_numbers.py} $\to$ \texttt{numbers.md} (single data-source mapping from body numbers to JSON fields and scripts); \texttt{params.csv} / \texttt{params\_table.md} (parameter specification exported by \texttt{export\_params.py} and cross-checked with \texttt{q1\_solve.py} and \texttt{q2\_lib.py}).
\end{itemize}

\noindent The directory names \texttt{Q1/}--\texttt{Q4/} stand for the Chinese original directories in the submitted supporting materials. The correspondence between body numbers and data files is as follows:
\begin{itemize}
\item Section~4.2.4, three area conventions $\to$ \texttt{Q1/results/q1\_area\_estimators.json} (\texttt{area\_estimators.py}).
\item Table~\ref{tab:4-4}, the clearance-robust variant geometry of Section~4.3.5, and the $95\%$ ellipse semi-axes $\to$ \texttt{Q2/results/q2b\_precision.json} (\texttt{q2b\_gamma.py}).
\item Section~4.3.7, the certified lower bound and gap $\to$ \texttt{Q2/results/q2b\_lower\_bound.json} (\texttt{q2b\_bound.py}).
\item Section~4.4.3 and Table~\ref{tab:4-5} $\to$ \texttt{Q3/results/q3\_summary.txt} (\texttt{q3\_solve.py}).
\item Section~4.5.3 and Table~\ref{tab:4-6}, scenario details $\to$ \texttt{Q4/results/q4\_scenario\_list.md} (\texttt{q4\_scenario\_list.py}).
\item Section~4.5.3, simple-repair baseline and the three-way comparison $\to$ \texttt{Q4/results/q4\_repair.json} (\texttt{q4\_repair.py}).
\item Section~6.2, $Q/Q_2$ comparison $\to$ \texttt{Q1/results/q1\_metric\_compare.json} (\texttt{q1\_metric\_compare.py}).
\item Tables D-1 and D-2, all ten pairs and the 2(2) timeline $\to$ \texttt{Q2/results/q2a\_results.json}, \texttt{Q2/results/q2b\_results.json}.
\end{itemize}

\noindent The complete source programs, run scripts, and result JSON files are provided in the supporting materials.

\subsection*{B.3 Full symbol table}
\begin{table}[htbp]
\centering
\caption{Full symbol list}
\label{tab:B-2}
{\small
\renewcommand{\arraystretch}{1.08}
\begin{tabularnb}{@{}p{38mm}p{92mm}p{30mm}@{}}
  \toprule
  Symbol & Meaning & Unit / value \\
  \midrule
  $O$ & Dispatch center & $(50,50)$~m \\
  $B_i$ & Building $i$ ($i=1,\dots,5$), planar range $[x_i^0,x_i^1]\times[y_i^0,y_i^1]$, height $h_i$ & m \\
  $T_a$; $\theta_a,R_a,H_a$; $T'_a$ & Alert $a$ and its target; azimuth (counterclockwise from east), horizontal distance, height; horizontal projection & $^\circ$, m, m; m \\
  $h_{\mathrm{cam}}$; $\Delta z_a$ & Camera height above ground; camera--target height difference $\Delta z_a=H_a-h_{\mathrm{cam}}$ & 2~m; m \\
  $p=(x,y)$; $r(p)$; $\varepsilon(p)$ & Center of the observation station; horizontal distance from camera to target projection; LOS elevation $\arctan(\Delta z_a/r)$ & m; m; $^\circ$ \\
  $[\varepsilon_{\min},\varepsilon_{\max}]$ & Feasible gimbal pitch range & $[-10^\circ,80^\circ]$ \\
  C1, C2, C3; $\Omega_a$ & Vehicle geometry, gimbal pitch, and LOS occlusion constraints; observable region $\Omega_a=\mathrm{C1}\cap\mathrm{C2}\cap\mathrm{C3}$ & ---; m$^2$ \\
  $Q(p)$; $p^*,r^*,\varepsilon^*,Q^*$ & Observation quality $1/r$; optimal station, horizontal distance, elevation, and quality in Problem~1 & m$^{-1}$; m, m, $^\circ$, m$^{-1}$ \\
  $m(p)$; $\varphi$ & Geometric safety margin; polar angle parameter of the optimal arc & m; $^\circ$ \\
  $V_j$; $V_j^0$ & Observation vehicle $j$ and its initial position ($j=1,\dots,5$) & m \\
  $R_{\mathrm{veh}}$; $v$; $\omega$ & Equivalent vehicle disk radius; travel speed; maximum gimbal single-axis angular rate & 1.5~m; 3~m/s; $30^\circ$/s \\
  $\alpha$ & Azimuth from station to target (pan command angle) & $^\circ$ \\
  $t_{\mathrm{tr}},t_{\mathrm{sl}},t_{\mathrm{obs}},t_{\mathrm{rs}}$ & Travel, slew, synchronized observation, and reset times & s ($t_{\mathrm{obs}}=5$) \\
  $d_i$; $d$ & Horizontal distance from station $i$ to target; single-station shorthand in Problem~3 (nominal 40~m) & m \\
  $\sigma_\alpha$; $\gamma$; $\beta$ & Single-station azimuth standard deviation; intersection angle of the two LOS; two-station baseline length & $1^\circ$ (rad in computation); $^\circ$; m \\
  $\Sigma$; $G$ & Covariance of the intersection-point position error; geometric dilution of precision $\sqrt{(d_1^2+d_2^2)/2}/|\sin\gamma|$ & m$^2$; m \\
  DRMS; $R_{95}$ & Root-mean-square position error $\sigma_\alpha G$; 95\% error-circle radius $2.4477\,\mathrm{DRMS}$ & m; m \\
  $\pi_j$; $p_{j,a}$; $M_a$ & Task order of vehicle $j$; station of vehicle $j$ for alert $a$; vehicle pair observing alert $a$ & ---; m; --- \\
  $t^{\mathrm{rd}}_{v,k}$; $t_a^{\mathrm{start}}$ & Ready time of the $k$-th task of vehicle $v$; observation start time of alert $a$ & s \\
  $P,\hat P,\delta$; $s$; $\sigma(s)=a+bs$ & True position, estimated position, and localization error; cumulative travel distance; single-axis localization standard deviation ($a,b\ge0$) & m; m; m, m/m \\
  $\varepsilon_0,\varepsilon_c$; $\Delta\psi$; $\sigma_\psi$; $s_{\max}$ & Gimbal bias/control error; difference between commanded and true azimuth ($\wrap$); combined standard deviation of $\Delta\psi$; maximum travel distance satisfying $P(|\Delta\psi|\le5^\circ)\ge0.95$ & $0.8^\circ/0.6^\circ$; $^\circ$; $^\circ$; m \\
  $t_p,t_{\mathrm{rec}}$; $N,k,\eta$; $t_{\mathrm{last}}$ & Rescheduling trigger/communication recovery time; total alerts, successful starts, $\eta=k/N$; latest observation start time (60~s if not started) & s; ---; s \\
  \bottomrule
\end{tabularnb}
}
\end{table}

% ---------------------------------------------------------------------
\section*{Appendix C: Figure and Table Inventory}
\addcontentsline{toc}{section}{Appendix C: Figure and Table Inventory}

The four problem figure sets contain 11 multi-panel combined figures (3 for Q1, 4 for Q2, 2 for Q3, and 2 for Q4), plus one framework figure in Chapter~2 (Figure~\ref{fig:framework}); each is delivered as an SVG vector with a PNG fallback.

\renewcommand{\thetable}{C-\arabic{table}}
\setcounter{table}{0}

\begin{table}[htbp]
\centering
\caption{Figure inventory}
\label{tab:C-1}
{\small
\renewcommand{\arraystretch}{1.08}
\begin{tabularnb}{@{}p{22mm}p{110mm}p{20mm}@{}}
  \toprule
  Figure & Content & Section \\
  \midrule
  Figure 2-1 & Framework of the four problems and the modeling chain & 2 \\
  Figure 4-1 & Observable region of T1 and building occlusion shadows & 4.2.4 \\
  Figure 4-2 & Optimal arc and margin--azimuth curve & 4.2.3 \\
  Figure 4-3 & Sensitivity of $r^*$ and the observable area & 4.2.4 \\
  Figure 4-4 & Observable region of T2 and 30~s reachable domains & 4.3.2 \\
  Figure 4-5 & Two-station accuracy versus $\gamma$, $d$, and baseline & 4.3.2 \\
  Figure 4-6 & Intersection error ellipses and the $95\%$ circle & 4.3.2 \\
  Figure 4-7 & Gantt chart and deployment routes of the time-optimal plan & 4.3.5 \\
  Figure 4-8 & Variance components and shares of $\sigma_\psi$ & 4.4.3 \\
  Figure 4-9 & Probability curve $P(|\Delta\psi|\le5^\circ)$ and the $(a,b)$ region & 4.4.3 \\
  Figure 4-10 & Two-stage response of $t_{\mathrm{last}}$ to $t_{\mathrm{rec}}$ & 4.5.3 \\
  Figure 4-11 & Rescheduling Gantt chart and the 74-scenario comparison & 4.5.3 \\
  \bottomrule
\end{tabularnb}
}
\end{table}

\begin{table}[htbp]
\centering
\caption{Table inventory}
\label{tab:C-2}
{\small
\renewcommand{\arraystretch}{1.08}
\begin{tabularnb}{@{}p{22mm}p{110mm}p{20mm}@{}}
  \toprule
  Table & Content & Section \\
  \midrule
  Table 2-1 & Building geometry and initial vehicle positions & 2.1 \\
  Table 2-2 & Compact algorithm routes (full table in B-1) & 2.6 \\
  Table 3-1 & Core symbols (full table in B-2) & 3.1 \\
  Table 4-1 & Main results of Problem~1 & 4.2.4 \\
  Table 4-2 & Two-station accuracy and station selection & 4.3.2 \\
  Table 4-3 & Scheduling plan versus the nearest-assignment baseline & 4.3.4 \\
  Table 4-4 & Intersection geometry of the timeliness plan & 4.3.6 \\
  Table 4-5 & Monte Carlo verification for Problem~3 & 4.4.3 \\
  Table 4-6 & Three-way comparison by scenario class for Problem~4 & 4.5.3 \\
  Table 4-7 & Comparison of representative scenarios & 4.5.3 \\
  Table 5-1 & Summary of validation for the four problems & 5 \\
  Table 6-1 & Objective time under the $v\times\omega$ full factorial & 6.2 \\
  Table B-1 & Full algorithm routes and selection rationale & Appendix B.1 \\
  Table B-2 & Full symbol list & Appendix B.3 \\
  Table D-1 & All ten vehicle pairs for T2 & Appendix D \\
  Table D-2 & Stations and per-task timeline of the 10.355~s plan & Appendix D \\
  \bottomrule
\end{tabularnb}
}
\end{table}

% ---------------------------------------------------------------------
\section*{Appendix D: Supplementary Tables}
\addcontentsline{toc}{section}{Appendix D: Supplementary Tables}

\renewcommand{\thetable}{D-\arabic{table}}
\setcounter{table}{0}

\begin{table}[htbp]
\centering
\caption{All ten vehicle pairs for T2 in 2(1)}
\label{tab:D-1}
{\footnotesize
\renewcommand{\arraystretch}{1.06}
\begin{tabularnb}{@{}p{14mm}p{26mm}p{14mm}p{16mm}p{16mm}p{24mm}p{26mm}@{}}
  \toprule
  Pair & $d_1/d_2$ (m) & $\gamma$ ($^\circ$) & $G$ (m) & DRMS (m) & $t_1/t_2$ (s) & Feasible within 30~s \\
  \midrule
  V1+V2 & 5.8188 / 11.2688 & 90.000 & 8.9678 & 0.1565 & 19.19 / 30.29 & No \\
  V1+V3 & 5.8188 / 5.8188 & 90.000 & 5.8188 & 0.1016 & 18.71 / 3.49 & Yes \\
  V1+V4 & 5.8188 / 5.8188 & 90.000 & 5.8188 & 0.1016 & 18.29 / 27.64 & Yes \\
  V1+V5 & 5.8188 / 5.8188 & 90.000 & 5.8188 & 0.1016 & 19.47 / 14.17 & Yes \\
  V2+V3 & 5.8188 / 11.2688 & 90.000 & 8.9678 & 0.1565 & 34.75 / 10.08 & No \\
  V2+V4 & 5.8188 / 11.2688 & 90.000 & 8.9678 & 0.1565 & 34.75 / 25.37 & No \\
  V2+V5 & 5.8188 / 11.2688 & 90.000 & 8.9678 & 0.1565 & 34.75 / 16.54 & No \\
  V3+V4 & 5.8188 / 5.8188 & 90.000 & 5.8188 & 0.1016 & 6.49 / 27.64 & Yes \\
  \textbf{V3+V5} & \textbf{5.8188 / 5.8188} & \textbf{90.000} & \textbf{5.8188} & \textbf{0.1016} & \textbf{5.11 / 13.61} & \textbf{Yes (selected)} \\
  V4+V5 & 5.8188 / 5.8188 & 90.000 & 5.8188 & 0.1016 & 28.95 / 15.42 & Yes \\
  \bottomrule
\end{tabularnb}
}
\par\smallskip
\begin{minipage}{\linewidth}
\footnotesize Note: $d_1,d_2$ are the station--target distances, $G$ and DRMS follow Section~4.3.1 with $\sigma=1^\circ=0.017453$~rad, and $t_1,t_2$ are the obstacle-avoiding travel times. The four V2 pairs share the same candidate station, hence their identical $G$. ``Feasible'' means both ready times are within the 30~s limit; among the six feasible pairs, V3+V5 has the smallest total travel time ($18.72$~s).
\end{minipage}
\end{table}

\begin{table}[htbp]
\centering
\caption{Stations and per-task timeline of the 10.355~s plan in 2(2)}
\label{tab:D-2}
{\footnotesize
\renewcommand{\arraystretch}{1.06}
\begin{tabularnb}{@{}p{12mm}p{10mm}p{28mm}p{13mm}p{13mm}p{12mm}p{16mm}p{14mm}p{14mm}@{}}
  \toprule
  Vehicle & Task & Station $(x,y)$ (m) & $\tau_{\mathrm{tr}}$ (s) & $\tau_{\mathrm{sl}}$ (s) & Start (s) & Observation end (s) & Reset end (s) & Clearance (m) \\
  \midrule
  V3 & T1 & (15.2004, 87.2331) & 0.57 & 1.03 & 1.607 & 6.607 & 7.641 & 5.581 \\
  V5 & T1 & (49.9000, 92.8000) & 0.07 & 1.53 & 1.607 & 6.607 & 8.140 & 7.080 \\
  V1 & T2 & (14.8000, 14.3000) & 0.12 & 3.12 & 6.994 & 11.994 & 15.117 & 7.295 \\
  V4 & T2 & (86.9000, 70.9000) & 1.03 & 5.96 & 6.994 & 11.994 & 17.953 & 0.067 \\
  V2 & T3 & (88.1000, 9.5000) & 0.17 & 5.75 & 10.355 & 15.355 & 21.100 & 10.897 \\
  V3 & T3 & (16.7656, 88.6590) & 0.71 & 2.01 & 10.355 & 15.355 & 17.363 & 0.001 \\
  \bottomrule
\end{tabularnb}
}
\par\smallskip
\begin{minipage}{\linewidth}
\footnotesize Note: the observation lasts 5~s and the reset time is the slew back to the east-facing attitude. The start time of an alert is the later of its two vehicles' ready times, so V1 waits from $3.24$~s to $6.994$~s for the longer slew of V4. V3 serves T1 and then T3 (the only reuse chain); T1 is already under way at $t_p=5$~s and is therefore locked in Problem~4. Clearance is the minimum distance from the station to the $1.5$~m-expanded building footprints; the two knife-edge values ($0.001$~m and $0.067$~m) are those discussed in Section~4.3.5.
\end{minipage}
\end{table}
